\section{Consequences for the quasi-polynomial programme}
\label{sec:research}

\subsection{What obstacle has been removed}

The new bound removes dependency depth from the cost of one prescribed
singleton system. A chain of a thousand definitions need not cause a
thousand complete rewrites of the relator grammar. All definitions and
their inverses can be composed together with a constant-factor bound on
the grammar size. The preserved dependency structure also admits a short
ordered witness whose validity can be checked without first constructing
all substituted words separately.

This improves the situation at a concrete point in the project. The old
forest policy was already capable of sharing long powers, but mixed signs
and additional dependencies excluded middle relations on a simple actual
diagram family. The new operation exposes those relations immediately and
avoids the resulting linear chain of scans and normalization handoffs.
The certified group-stage measurements agree with that explanation.

The eager-policy regression shows a different obstacle. A smaller
description after an exact quotient need not be a better input to an
existing heuristic. Search quality depends on which relations remain
visible, which letters are repeated, and what subsequent transformations
become cheap. Grammar size, the number of live generators, expanded length,
and the cost of later search are different quantities. The final terminal
guard avoids needing to predict the latter for this release.

\subsection{A useful amortized planning bound}

The delivered terminal probe can run at many search states without using
the general batch at any of them. Its additional discovery cost is therefore
worth accounting for separately.

\begin{proposition}[Cached planning across a trace]
\label{prop:amortized-planning}
Suppose a search trace invokes the singleton planner $h$ times, retains
$m$ relator slots, and never has more than $r_0$ live generators. Let $Q$
be the number of distinct arena nodes whose immutable maximum-and-count
summary is first encountered across these invocations. The planner's
charged structural work is
\[
 O\!\left(Q+h\bigl(m+r_0\log(r_0+1)+1\bigr)\right),
\]
apart from the ordinary bit cost of keys and label comparisons.
\end{proposition}

\begin{proof}
Each immutable node's summary is constructed once. Its binary children
give only a bounded number of stack events, so all fresh metadata costs
$O(Q)$. A root already in the cache stops its descent immediately.
Current slots, candidate membership, and distinct-child selection cost
$O(m+r_0)$ per invocation. The selected list has at most $r_0$ entries;
the explicit sorting allowance costs $O(r_0\log(r_0+1))$.
Summing over the invocations gives the bound.
\end{proof}

This result avoids an unnecessary factor $hM$ for repeatedly rescanning
unchanged shared subgrammars. It does not bound $Q$ or $h$ for the complete
recognizer. Whitehead moves, overlap operations, and normalization can
create additional nodes and additional states; their costs remain part of
the larger research problem. It also explains why an exact early probe can
still add a modest measured overhead when it seldom fires.

\subsection{What a general quasi-polynomial theorem still requires}

For the algebraic route, a complete proof needs a quantitative exposure
result: every unknot in the chosen input model must reach a useful terminal
or another complete decision method after controlled preparation. Such a
result would have to bound the number of states visited, the size of their
shared representations, the work of normalization and relator discovery,
and certificate replay. Efficient primitives establish none of these
global existence and scheduling statements by themselves.

The raw phase theorem in Section~\ref{alg:complexity} can be used once its
hypotheses are actually supplied. In particular, $d=O((\log n)^2)$ pure
singleton phases from polynomial-size data yield an
$n^{O(\log n)}$ raw-contraction bound. With $d=O(\log n)$ they yield a
polynomial raw-contraction bound. Interleaved normalization or unrelated
moves require their own bounds; they cannot be placed inside the same
constant-factor recurrence without proof. The delivered guarded search
uses at most one singleton batch in a completed certificate, but may
perform many preceding search steps.

For the geometric route suggested by Lackenby's programme, the key
obligations remain constructive surface selection, encoded cutting and
attachment data, effective parallelity-bundle operations, controlled
hierarchy depth, and the total bit cost of all intermediate objects
\cite{Lackenby2026Hierarchy}. Agol--Hass--Thurston orbit counting is a
foundational example of obtaining polynomial work in a compressed normal
surface description without expanding its pieces \cite{AgolHassThurston2006}.
Lackenby's fast curve algorithms offer related boundary-surface primitives
\cite{Lackenby2026Curves}. These results are relevant tools, but they do
not imply that the missing hierarchy can be selected within the target
budget. This article does not identify its presentation DAG with a
three-manifold hierarchy.

\subsection{Further questions with concrete success criteria}

\paragraph{1. Better orders with bounded discovery cost.}
The fixed numeric order misses valid batches, even on relabelled stars.
Can one compute a small deterministic portfolio of orders that exposes
substantially more terminal batches on actual Wirtinger states? A useful
result would combine a restricted structural theorem with complete-call
measurements that charge every unsuccessful order. The NP-completeness
result for unrestricted positive presentations rules out treating maximum
batch selection as an automatically easy preprocessing problem. It does
not rule out efficient algorithms for the planar source class or useful
parameterized algorithms.

\paragraph{2. A structural characterization of terminal exposure.}
Which knot diagrams have a Wirtinger presentation admitting a singleton
DAG to one surviving generator, before any further word transformation?
Existence is invariant under a pure generator renaming, because the witness
can be transported with its order; the numeric planner's success need not
be invariant. How does exposure change under a different presentation
construction or Reidemeister simplification? A characterization should
distinguish existence of a witness from success of one chosen order. It should
also specify whether the input is the raw four-letter presentation or
the initially freely and cyclically reduced version.

\paragraph{3. Quantitative exposure after controlled transformations.}
Can a diagram class broader than descending or visibly stabilized unknots
be shown to reach a terminal singleton system after polynomial or
polylogarithmically many bounded transformations? The desired theorem
must charge the construction of those transformations, not merely show
that some sequence exists. This would be a direct route from the local
result to a stronger restricted recognition theorem.

\paragraph{4. Partial batches with a proved handoff guarantee.}
The Gordian experiment supplies a concrete rejection test for future
partial-batch policies. A replacement policy should retain the positive
coverage and compare the whole checked stage with the guarded version.
Mathematically, seek a potential that controls both representation size
and the next search operation, or a reversible scheduling scheme that
preserves the intact predecessor state with a proved total-cost bound.
Simply choosing the largest available batch is not justified by the
present experiments.

\paragraph{5. Raw closure with persistent substitutions.}
Can several raw singleton phases be represented with an additive
description bound rather than the conservative constant-factor recurrence?
The tempting proof by unioning original donor-pivot pairs is invalid:
after eliminating $x$ using $xg$, the original slot $x$ exposes the newly
introduced singleton $g$. A successful construction must track that
provenance, empty images, and newly exposed occurrences. It should also
bound the discovery cost of the closure itself.

\paragraph{6. Memory reclamation consistent with source replay.}
The current arena retains nodes that become unreachable. The theorem
bounds new allocations relative to the reachable source, and explicitly
adds pre-existing unreachable storage. Can old nodes be compacted with a
verified remapping that preserves caches, certificate provenance, and
the cost of subsequent comparisons? A useful implementation should report
live nodes, allocated nodes, auxiliary metadata bits, and wall time
separately.

\paragraph{7. General primitive and proper-power donors.}
The new operation is strongest when the donor itself has a literal
singleton. Existing primitive-power certificates can infer additional
definitions under source-established torsion-freeness. Can a larger
mixed batch combine these two kinds of evidence while preserving linear
context construction and independent replay? Non-unit rank-two primitive
relations and cyclic dependency components require their own quotient
analysis; they cannot simply be admitted as additional DAG edges.

\paragraph{8. Endpoints with a small remaining rank.}
Rank one has a particularly cheap and complete exponent criterion.
For higher rank, identify restricted residual presentations admitting
equally explicit complete checks, then allow a batch only when it reaches
one of those checked endpoints. A theorem parameterized by the actual
residual structure would be more informative than one parameterized only
by the number of remaining generator names.

\paragraph{9. Certificates as separate reusable mathematical objects.}
The ordered witness describes a presentation equivalence independently of
the search that found it. A Lean formalization of the quotient theorem,
the ranked-summary checker, and the signed grammar evaluator would make
that separation precise. It would still need an independent account of
diagram-to-presentation reconstruction and the classical topological
endpoint before becoming a formal unknot recognizer.

\paragraph{10. Representative benchmarks and crossover prediction.}
The actual source-family speedups are large, whereas the ordinary corpus
shows overhead. A next benchmark should contain more diagrams that survive
the early filters, preserve source and verdict provenance, and compare
features under both generous and production resource allowances. Prediction
of when to invoke the terminal probe should be assessed using withheld
cases. Measurements should retain all nondecisions and A/A controls and
should distinguish a faster group stage from a faster complete recognizer.

\subsection{Recommended next research priority}

The most direct next mathematical target is a restricted exposure theorem
coupled to a small order portfolio. It would address a condition that the
current implementation can already exploit and check. The most direct
engineering target is a cheap, evidence-based invocation policy for the
guarded terminal probe. Both can use the delivered corpus, the eager
Gordian counterexample, and the independent checker as fixed evaluation
tools. A general quasi-polynomial claim should wait for the corresponding
global size and work bounds.
